\perchtable{tab-status}
\perchfigure{fig-recovery}

\paragraph{Sources.} \ifdefined\anonymous The revisions of Perch, rule-core,
soroban-sim and the Guess the Number fixture are withheld for review and given in
the artifact.\else This paper describes Perch at revision
\facthash{perch-pin}~\selfcite{perch-src} and rule-core at
\facthash{rule-core-pin}~\selfcite{rule-core}. The reference run uses soroban-sim
at \facthash{soroban-sim-pin} and the Guess the Number fixture at
\facthash{guess-the-number-pin}~\selfcite{guess-the-number}.\fi

\paragraph{Proofs.} The proofs reproduce from rule-core with Charon
\fact{charon-rev}, Aeneas \fact{aeneas-rev} and Lean \fact{lean-rule-core}:
\code{tools/extract.sh -{}-check} regenerates the translation and fails on any
difference, and \code{lake build} in \code{formal/} checks the theorems;
rule-core's README gives the full commands. Perch's canonical-encoding package
uses Lean \fact{lean-perch-canon}. rule-core's CI also checks a ledger of
\fact{claims-total} claims (\fact{claims-proven} proven, \fact{claims-tested}
tested, \fact{claims-assumed} assumed, \fact{claims-unproven} unproven).

\paragraph{Reference run.} Setup loads protocol-\fact{sim-protocol} cost and
configuration entries from a testnet fixture (ledger \fact{fixture-ledger};
testnet now runs protocol \fact{testnet-protocol}) and places the compiler and the
Perch \fact{perch-deployed-account-version} policy contracts at the addresses the
account derives from its pins; a \fact{registry-stub-bytes}-byte registry stub
resolves the verifier name \code{\fact{sim-verifier-name}} offline, so the
production registry's governance is not tested. Refusals caught in enforcing simulation are not submitted,
so an unchanged fingerprint says nothing about fees for failed, submitted
transactions. A native program and a WASI reactor run the same driver over the same pinned bytes and agree on every
returned field except wall-clock time. The HTML edition embeds a source bundle
with the Rust JSON-command driver, its native entry point, a
\fact{wasi-imports-count}-function WASI shim, the browser worker, dependency
locks, the contract build script, the reference command sequence and a
comparison checker; \code{projects/perch-browser/README.md} in the bundle gives
the pinned build commands. Like Perch's own release build, the build rewrites the
account's infrastructure pin files with the hashes of the contracts built in the
same run, so the embedded account (\fact{embedded-account-short}) differs from
the testnet preview's (\fact{preview-account-wasm-short}). The account still pins
the Perch \fact{perch-deployed-account-version} policy contracts and reports them
through \code{infra()}, but never calls them. The demo's driver maps the names
this paper uses to the pinned exports of Table~\ref{tab:tab-api} at its call
boundary. An executable walkthrough of the task in
Table~\ref{tab:tab-walkthrough}, in both designs, accompanies the HTML edition.

\paragraph{This paper.} \code{python3 docs/perch-paper/build/paper.py} renders
both editions, the HTML edition and this PDF, from shared sources, and
\code{docs/perch-paper/latex/build.sh} builds this PDF with a pinned Tectonic.

\perchtable{tab-run}

\perchtable{tab-api}

\perchtable{tab-artifacts}
\perchtable{tab-ending}
